%% ---- Supplementary Note [note:graphs]: the gene--gene graphs used as attention priors ----
%% Every number here is emitted by experiments/010-kuzmin-tmi/scripts/graph_statistics.py (graph
%% sizes, overlap, STRING releases; also writes tab-graphs.tex + tab-string-versions.tex) and
%% experiments/005-kuzmin2018-tmi/scripts/dango_string_version_sweep.py (wandb pull of the DANGO
%% STRING-release sweep; writes tab-dango-string-versions.tex). The figures are composed by
%% experiments/010-kuzmin-tmi/scripts/compose_graph_si_figures.py. Rerun those, not this file.
\clearpage
\suppnote{note:graphs}{gene--gene graphs used as attention priors\secstatus{todo}}

Nine gene--gene graphs regularize the cell encoder's attention in the trigenic experiment,
one head per graph, and every one of them enters the model only through the soft prior
$\Omega_k$ of Eq.~\ref{eq:graphreg}: none is used for message passing and none is imposed as a
hard mask. \supptab{tab:databases} and \suppfig{fig:graph-priors} report how large the graphs
are and how much they overlap, \supptab{tab:string-versions} and \suppfig{fig:dango-string}a
how the STRING channels have drifted across releases, and \supptab{tab:dango-string-versions}
and \suppfig{fig:dango-string}b whether that drift moves the DANGO baseline.

\notesec{Sources} Two graphs are curated from the \emph{Saccharomyces} Genome Database
(SGD). \emph{Physical} holds the protein--protein interactions SGD records per gene.
\emph{Regulatory} holds SGD's regulation annotations as directed regulator$\to$target edges;
of its 39,636 records, 38,677 are transcriptional and 37,890 come from high-throughput
assays (24,199 from chromatin immunoprecipitation), so 511 regulator genes reach 6,546 target
genes and the graph touches 6,582 of the 6,607 genes in the reference. \emph{TFLink} is the
transcription-factor--target compendium of the same name (v1.0), also directed. The remaining
six are the evidence channels STRING v12.0 scores separately for \emph{S.\ cerevisiae}: gene
neighborhood, gene fusion, phylogenetic co-occurrence, co-expression, experimental
(protein-interaction assays), and curated database. A channel edge is any gene pair with a
nonzero channel score; no score threshold is applied, matching the cut DANGO trains on. Every
graph is restricted to the S288C reference gene set of 6,607 genes, and the directed graphs
are treated as undirected gene pairs wherever graphs are compared below.

\notesec{Size and coverage} The graphs differ by a factor of three in the genes they reach
(2,191 for neighborhood, 6,474 for co-expression) and by two orders of magnitude in edges
(11,085 for co-occurrence, 996,199 for co-expression), with mean degree from 7.7 (fusion) to
308 (co-expression) (\suppfig{fig:graph-priors}a; \supptab{tab:databases}). Their degree
distributions differ in shape as well as scale (\suppfig{fig:graph-priors}b): fusion and
co-occurrence have no gene above degree 100, the two dense STRING channels fall off sharply
past degree $10^{3}$, and TFLink carries the heaviest tail, with transcription factors
reaching several thousand targets. Summed over the nine graphs the edges name 2,437,196 gene
pairs, of which 1,515,940 are distinct, 6.9\% of the 21.8 million pairs the vocabulary
allows.

\notesec{Overlap} The graphs are complementary rather than redundant. The Jaccard index of
edge sets exceeds 0.12 for a single pair, co-expression against experimental (0.44)
(\suppfig{fig:graph-priors}c). Containment is strongly asymmetric
(\suppfig{fig:graph-priors}d): 76\% of physical edges lie inside the experimental channel and
51\% inside co-expression, 82\% of neighborhood edges lie inside co-expression, and 43\% of
regulatory edges lie inside TFLink, whereas only 8\% of TFLink edges are regulatory. The two
dense channels therefore absorb much of the sparse ones without the converse holding. Counting
support per pair (\suppfig{fig:graph-priors}e), 52\% of the distinct pairs appear in exactly
one graph, 37\% in two, and none in more than seven. TFLink (80\%) and regulatory (51\%)
contribute the largest shares of edges found nowhere else; database (8\%) and neighborhood
(10\%) the smallest (\supptab{tab:databases}, Unique). How much each graph contributes to
held-out accuracy is not measured here.

\notesec{STRING releases} DANGO was built on STRING v9.1 and set its per-channel weights from
the change to v11.0; the CGT experiment uses v12.0. Across the three releases the channels
grow by 2.2$\times$ (database) to 8.7$\times$ (fusion) in edges
(\supptab{tab:string-versions}; \suppfig{fig:dango-string}a), and the releases are not nested:
only 42\% (fusion) to 72\% (neighborhood) of a v9.1 channel's pairs persist in v11.0, the
Jaccard index between consecutive releases is 0.12--0.35, and between v9.1 and v12.0 it is
0.03--0.19. A model trained on one release therefore sees a substantially different prior from
one trained on the next.

\notesec{DANGO replication} The TorchCell reimplementation of DANGO was trained on the
Kuzmin~2018 trigenic interactions~\cite{kuzminSystematicAnalysisComplex2018a} with the six
channels taken from v9.1, v11.0, or v12.0, under each of three schedules for handing over
from the pretraining reconstruction loss to the interaction loss (pretrain then main; linear
to uniform; linear to flipped). The statistic is the maximum over epochs of validation
Pearson~$r$, the checkpoint-selection rule, an upward-biased order statistic whose epoch count
is reported beside it. Across the 19 runs the best validation Pearson lies between 0.415 and
0.427, and no release or schedule separates from the others beyond the run-to-run spread
(\supptab{tab:dango-string-versions}; \suppfig{fig:dango-string}b; SEM at most 0.003 where
$n>1$). Swapping v9.1 for v12.0 neither helps nor hurts the baseline, even though the v12.0
channels carry three to nine times as many edges. These are validation maxima on the DANGO
split and are not the test-split baseline of Fig.~\ref{fig:ggi}.

%% Figure source: notes/assets/drawio/FigS-graph-attention-priors.drawio (written by
%%   experiments/010-kuzmin-tmi/scripts/compose_graph_si_figures.py from the panels of
%%   experiments/010-kuzmin-tmi/scripts/graph_statistics.py)
%%   -> figures/FigS-graph-attention-priors.pdf (auto-exported by `make fig`/`make paper`).
\begin{figure*}[t]
\tcfig{figures/FigS-graph-attention-priors.pdf}{\textbf{Gene--gene graphs used as attention
priors in the trigenic experiment.} (a)~Genes with at least one edge (dotted line, the
6,607-gene reference), distinct gene pairs, and the fraction of each graph's pairs found in no
other graph. (b)~Complementary cumulative degree distribution over undirected gene pairs.
(c)~Jaccard index of edge sets. (d)~Containment: the fraction of the row graph's pairs present
in the column graph. (e)~Distinct gene pairs supported by exactly $n$ of the nine graphs. Source:
\texttt{experiments/010-kuzmin-tmi/scripts/graph\_statistics.py}.}{fig:graph-priors}
\end{figure*}

%% Figure source: notes/assets/drawio/FigS-dango-string-versions.drawio (same composer; panel b
%%   from experiments/005-kuzmin2018-tmi/scripts/dango_string_version_sweep.py)
%%   -> figures/FigS-dango-string-versions.pdf (auto-exported by `make fig`/`make paper`).
\begin{figure*}[t]
\tcfig{figures/FigS-dango-string-versions.pdf}{\textbf{STRING releases and the DANGO
replication.} (a)~Edges per STRING channel in v9.1, v11.0, and v12.0 (top) and the Jaccard
index between releases of the same channel (bottom). (b)~DANGO replication on the Kuzmin~2018
trigenic interactions: best validation Pearson~$r$ (maximum over epochs) by loss schedule and
STRING release; bar, mean over runs; whisker, SEM where $n>1$; circle, one run (19 runs).
Sources: \texttt{graph\_statistics.py} (a) and
\texttt{experiments/005-kuzmin2018-tmi/scripts/dango\_string\_version\_sweep.py} (b).}{fig:dango-string}
\end{figure*}

%% SOURCE: experiments/010-kuzmin-tmi/scripts/graph_statistics.py -- AUTO-GENERATED, do not hand-edit; rerun the script.
%% Node = gene with at least one edge; Edges = native count (directed for regulatory and
%% TFLink); Pairs = distinct unordered gene pairs; Unique = pairs found in no other graph.
\begin{table}[t]
\centering
\footnotesize
\caption{Gene--gene graphs that serve as attention priors in the trigenic experiment,
over the shared S288C vocabulary of 6,607 genes. \textbf{Physical} and \textbf{Regulatory}
are the curated SGD graphs; TFLink and the six STRING~v12.0 evidence channels add the
remaining rows. \emph{Edges} is the native count (directed for Regulatory and TFLink),
\emph{Pairs} the distinct unordered gene pairs used for overlap, \emph{Unique} the share of
those pairs found in no other graph. The nine graphs together cover
1,515,940 distinct gene pairs; relative to the SGD-native baseline (Physical\,$+$\,Regulatory),
the sum over graphs is 3.39$\times$ in nodes and
13.63$\times$ in edges.
(STRING~v9.1 and v11.0, used in the DANGO replication, are in \supptab{tab:string-versions}.)}
\label{tab:databases}
\begin{tabular}{@{}l r r r r r@{}}
\toprule
\textbf{Graph} & \textbf{Nodes} & \textbf{Edges} & \textbf{Pairs} & \textbf{Mean degree} & \textbf{Unique}\\
\midrule
Physical (SGD) & 5,721 & 139,463 & 137,604 & 48.1 & 23\%\\
Regulatory (SGD) & 6,582 & 39,636 & 39,546 & 12.0 & 51\%\\
TFLink & 5,074 & 200,801 & 199,551 & 78.7 & 80\%\\
STRING neighborhood & 2,191 & 146,713 & 146,713 & 133.9 & 10\%\\
STRING fusion & 3,081 & 11,787 & 11,787 & 7.7 & 14\%\\
STRING co-occurrence & 2,601 & 11,085 & 11,085 & 8.5 & 20\%\\
STRING co-expression & 6,474 & 996,199 & 996,199 & 307.8 & 35\%\\
STRING experimental & 6,016 & 822,094 & 822,094 & 273.3 & 25\%\\
STRING database & 4,027 & 72,617 & 72,617 & 36.1 & 8\%\\
\midrule
Sum (all graphs) & 41,767 & 2,440,395 & 2,437,196 & & \\
Union of pairs & & & 1,515,940 & & \\
Sum (Physical, Regulatory) & 12,303 & 179,099 & 177,150 & & \\
\bottomrule
\end{tabular}
\end{table}

%% SOURCE: experiments/010-kuzmin-tmi/scripts/graph_statistics.py -- AUTO-GENERATED, do not hand-edit; rerun the script.
\begin{table}[t]
\centering
\footnotesize
\caption{STRING evidence channels for \emph{S.\ cerevisiae} across the three releases
integrated into TorchCell, over the shared 6,607-gene vocabulary. Nodes are genes with at
least one edge; edges are distinct unordered gene pairs with a nonzero channel score.
DANGO was published on v9.1 with a v9.1$\to$v11.0 comparison; the CGT trigenic
experiment uses v12.0.}
\label{tab:string-versions}
\begin{tabular}{@{}l rrr rrr@{}}
\toprule
 & \multicolumn{3}{c}{\textbf{Nodes}} & \multicolumn{3}{c}{\textbf{Edges}}\\
\cmidrule(lr){2-4}\cmidrule(lr){5-7}
\textbf{Channel} & v9.1 & v11.0 & v12.0 & v9.1 & v11.0 & v12.0\\
\midrule
neighborhood & 2,170 & 2,115 & 2,191 & 45,601 & 120,965 & 146,713\\
fusion & 1,187 & 2,338 & 3,081 & 1,358 & 3,914 & 11,787\\
co-occurrence & 1,268 & 1,307 & 2,601 & 2,659 & 4,668 & 11,085\\
co-expression & 5,796 & 5,654 & 6,474 & 313,688 & 554,833 & 996,199\\
experimental & 6,079 & 5,964 & 6,016 & 219,450 & 392,772 & 822,094\\
database & 2,712 & 3,081 & 4,027 & 33,486 & 46,816 & 72,617\\
\midrule
Sum & 19,212 & 20,459 & 24,390 & 616,242 & 1,123,968 & 2,060,495\\
\bottomrule
\end{tabular}
\end{table}

%% SOURCE: experiments/005-kuzmin2018-tmi/scripts/dango_string_version_sweep.py -- AUTO-GENERATED from wandb zhao-group/torchcell_005-kuzmin2018-tmi_dango; do not hand-edit.
\begin{table}[t]
\centering
\footnotesize
\caption{DANGO replication in TorchCell on the Kuzmin 2018 trigenic interactions: best
validation Pearson $r$ (maximum over epochs of the logged validation Pearson, the
checkpoint-selection rule) by STRING release and loss schedule. Mean $\pm$ SEM over $n$
runs; a single run reports its value with no whisker. Epochs is the range of epochs logged
across those runs.}
\label{tab:dango-string-versions}
\begin{tabular}{@{}l l r l l@{}}
\toprule
\textbf{STRING} & \textbf{Loss schedule} & $n$ & \textbf{Val.\ Pearson $r$} & \textbf{Epochs}\\
\midrule
v9.1 & pretrain then main & 1 & 0.420 & 1000\\
v9.1 & linear to uniform & 1 & 0.422 & 575\\
v9.1 & linear to flipped & 2 & 0.422 $\pm$ 0.001 & 500--540\\
v11.0 & pretrain then main & 1 & 0.419 & 988\\
v11.0 & linear to uniform & 2 & 0.424 $\pm$ 0.003 & 443--468\\
v11.0 & linear to flipped & 2 & 0.422 $\pm$ 0.001 & 500\\
v12.0 & pretrain then main & 4 & 0.423 $\pm$ 0.002 & 500\\
v12.0 & linear to uniform & 5 & 0.420 $\pm$ 0.001 & 500\\
v12.0 & linear to flipped & 1 & 0.422 & 500\\
\bottomrule
\end{tabular}
\end{table}

